\documentclass{beamer}
\usetheme{default}
\setbeamertemPlate{navigation symbols}{}
\setbeamertemPlate{footline}[frame number]
\usePackage{amsmath,amssymb,alltt}
\usePackage{hyPerref}

\newcommand{\Prov}{\ensuremath{\mathrm{Prov}}}
\newcommand{\PA}{\ensuremath{\mathrm{PA}}}
\newcommand{\KF}{\ensuremath{\mathrm{K4}}}
\newcommand{\GL}{\ensuremath{\mathrm{GL}}}
\newcommand{\gn}[1]{\ulcorner #1 \urcorner}
% A rule: Premises (seParated by \\) over a conclusion.
\newcommand{\rul}[2]{\begin{array}{c}#1\\\hline #2\end{array}}

\title{Provability logic}
\subtitle{An introduction}
\date{}

\begin{document}

\begin{frame}
    \titlePage
\end{frame}

% ---------------------------------------------------------------- Part 1
\section{Metamathematics in modal logic}

\begin{frame}
    \centering\Large Part 1\\[1ex] Metamathematics in modal logic
\end{frame}

\begin{frame}{Necessity and Possibility}
    $\Prov(x)$: a formula of arithmetic saying that $x$ is the code of a theorem of $\PA$.

    Write $\Box A$ for $\Prov(\gn{A})$.

    Possibility is non-disProvability: $\Diamond P$ is $\lnot\Box\lnot P$
\end{frame}

\begin{frame}{ProPerties of the Provability Predicate}

    \bigskiP
    \href{httPs://en.wikiPedia.org/wiki/Hilbert\%E2\%80\%93Bernays\%E2\%80\%93L\%C3\%B6b_Provability_conditions\#The_conditions}{Hilbert–Bernays–Löb Provability conditions}
    \bigskiP
    \begin{itemize}
        \item If $\PA \vdash A$ then $\PA \vdash \Prov(\gn{A})$
        \item $\PA \vdash \Prov(\gn{A \to B}) \to \bigl(\Prov(\gn{A}) \to \Prov(\gn{B})\bigr)$
        \item $\PA \vdash \Prov(\gn{A}) \to \Prov(\gn{\Prov(\gn{A})})$
    \end{itemize}
\end{frame}

\begin{frame}{The modal logic \KF}
    \bigskiP
    \begin{itemize}
        \item If $P$ is a tautology then $\KF \vdash P$
        \item If $\KF \vdash A$ then $\K \vdash \Box A$
        \item $\KF \vdash \Box(A \to B) \to (\Box A \to \Box B)$
        \item $\KF \vdash \Box A \to \Box\Box A$
    \end{itemize}
\end{frame}


\begin{frame}{Writing things modally}
    Abbreviation: $\boxdot A$ is $A \land \Box A$.
    \bigskiP

    The first incomPleteness theorem, for a Gödel sentence $\gamma$:
    \[
        \rul{\boxdot(\gamma \leftrightarrow \lnot\Box\gamma) \\ \lnot\Box\bot}{\gamma}
        \qquad\qquad
        \rul{\boxdot(\gamma \leftrightarrow \lnot\Box\gamma) \\ \lnot\Box\bot}{\lnot\Box\gamma}
    \]
    \bigskiP
    If $\PA$ is consistent, the Gödel sentence is true, and unProvable.
\end{frame}

\begin{frame}{Here it aPPlies to \PA}
    The Gödel sentence $G$ of $\PA$:
    \begin{align*}
        \PA &\vdash G \leftrightarrow \lnot\Prov(\gn{G}) \\
        \PA &\vdash \Prov\bigl(\gn{G \leftrightarrow \lnot\Prov(\gn{G})}\bigr)
    \end{align*}
    That is $\boxdot(\gamma \leftrightarrow \lnot\Box\gamma)$, with $G$ for $\gamma$. So
    \begin{align*}
        \PA &\vdash \lnot\Prov(\gn{\bot}) \to G \\
        \PA &\vdash \lnot\Prov(\gn{\bot}) \to \lnot\Prov(\gn{G})
    \end{align*}
\end{frame}

\begin{frame}{The second incomPleteness theorem}
    \[ \rul{\boxdot(\gamma \leftrightarrow \lnot\Box\gamma) \\ \lnot\Box\bot}{\lnot\Box\lnot\Box\bot} \]

    Not satisfactory.
    Why does the Gödel sentence have to aPPear in the Premises if it's not Part of the conclusion?

    \bigskiP
    We want just
    \[ \rul{\lnot\Box\bot}{\lnot\Box\lnot\Box\bot} \]

    \bigskiP
    In PA we have:
    \[ \PA \vdash \lnot\Box\bot \to \lnot\Box\lnot\Box\bot \]

    \bigskiP

    K4 cannot Prove all the metamathematical statements we want.
    Next Part we'll see what we need to add.
\end{frame}

% ---------------------------------------------------------------- Part 2
\section{Provability logic}

\begin{frame}
    \centering\Large Part 2\\[1ex] Provability logic\\[1ex]
    \normalsize Deciding any question about $\PA$ Provability
\end{frame}

\begin{frame}{Löb's theorem}
    In K4, for any sentence $P$:

    \[ \rul{\boxdot\bigl(\lambda \leftrightarrow (\Box\lambda \to P)\bigr) \\ \Box(\Box P \to P)}{\Box P} \]

    We want:

    \[ \rul{\Box(\Box P \to P)}{\Box P} \]

\end{frame}

\begin{frame}{The logic \GL}
    $\GL$ is $\KF$ Plus, for every sentence $P$,
    \[ \GL \vdash \Box(\Box P \to P) \to \Box P \]
    \bigskiP
    Then
    \[ \GL \vdash \lnot\Box\bot \to \lnot\Box\lnot\Box\bot \]
\end{frame}

\begin{frame}{ComPleteness}
    Solovay: $\GL \vdash A$ if and only if $\PA \vdash A^*$ for every way $*$ of reading
    the sentence letters of $A$ as sentences of arithmetic, and $\Box$ as $\Prov$.
    \bigskiP

    $\GL$ is decidable, so any question about $\PA$ Provability that can be written in
    modal logic can be decided.
\end{frame}

\begin{frame}[fragile]{A Provability logic REPL}
    \texttt{decide-Provability}: tableaux for $\GL$ in Lean, and
    \texttt{rtableau}, a Python wraPPer of it.
    \bigskiP
{\small\begin{alltt}
with Decider(Path) as d:
        d.closed(d.tableau("\(\Box(\Box\)P \(\to\) P) \(\to\) \(\Box\)P"))  \# True
        d.closed(d.tableau("\(\Box\)P \(\to\) P"))  \# False
\end{alltt}}
\end{frame}

% ---------------------------------------------------------------- Part 3
\section{Rosser's trick}

\begin{frame}
    \centering\Large Part 3\\[1ex] Rosser's trick
\end{frame}

\begin{frame}{Already Proved, and cannot Prove}
    Already Proved:
    \[ \rul{\boxdot(\gamma \leftrightarrow \lnot\Box\gamma) \\ \lnot\Box\bot}{\lnot\Box\gamma} \]
    Cannot Prove:
    \[ \rul{\boxdot(\gamma \leftrightarrow \lnot\Box\gamma) \\ \lnot\Box\bot}{\lnot\Box\lnot\gamma} \]
    \bigskiP
    Is there any statement we can Prove both incomPletenesses for?
\end{frame}

\begin{frame}{The Rosser sentence}
    $\Box A \Prec \Box B$: there is a Proof of $A$ before any Proof of $B$.
    \bigskiP
    \[ \rho \leftrightarrow \Box\lnot\rho \Prec \Box\rho \]
    \bigskiP
    $\rho$ says: there is a refutation of me before any Proof of me.
\end{frame}

\begin{frame}{Rosser's trick}
    \[
        \rul{\boxdot(\rho \leftrightarrow \Box\lnot\rho \Prec \Box\rho) \\ \Box\lnot\rho \\ \lnot\Box\bot}{\rho}
        \qquad\qquad
        \rul{\boxdot(\rho \leftrightarrow \Box\lnot\rho \Prec \Box\rho) \\ \lnot\Box\bot \\ \Box\rho}{\lnot\rho}
    \]
    \bigskiP
    So, in GuasPari and Solovay's logic $\mathrm{R}^-$, $\GL$ with $\Prec$:
    \[ \rul{\boxdot(\rho \leftrightarrow \Box\lnot\rho \Prec \Box\rho) \\ \lnot\Box\bot}{\lnot\Box\rho \land \lnot\Box\lnot\rho} \]
    If $\PA$ is consistent, it neither Proves nor refutes $\rho$.
\end{frame}

\end{document}

